\documentclass[11pt]{article}
\title{The released frontier in a selective route replacement}
\begin{document}
\maketitle

VDA 5050 distinguishes the released base from the mutable horizon.
The old released prefix may continue while the horizon is changed.
A release that has already been sent but has not been acknowledged
cannot be treated as part of the definitely unreleased horizon.

\paragraph{Acceptance witness.}
An initial order has released decision point \(A\), with \(B\)
unreleased in its horizon. Update 1 extends the base from \(A\) to
\(B\). Update 2 has the same order identity, starts at \(A\), and
releases the alternative \(X\), omitting \(B\). Both graph messages
validate in isolation. In the selected client code:
\begin{itemize}
\item If update 2 arrives first, it stitches at \(A\). Update 1 is
then rejected because the update number is not strictly increasing.
\item If update 1 arrives first, the released anchor is now \(B\).
Update 2 fails to stitch at \(A\) despite its greater update number.
A different update 2 anchored at \(B\) passes but leaves the earlier
release intact.
\end{itemize}
The replay checks these exact verdicts using the upstream validator
body and substitute field-compatible structs and error plumbing.
Node attributes outside the fixture are held equal at stitch points.
Delivery order is a test input: this does not assert that one healthy
MQTT session reorders messages. The distinction also arises if release
delivery is uncertain and the update is the next processed message.

\paragraph{Control consequence.}
If the old release has never been sent and cannot be issued by another
publisher, the retained prefix can finish while the fleet computes
the replacement. There is no release of \(B\) to rescind. After
dispatch of update 1, absence of its acknowledgement alone is
consistent with both acceptance histories. Sending update 2 at \(A\)
does not certify that \(B\) is excluded. Within the one-order model,
evidence that the robot accepted the specific update 2 at \(A\)
does resolve the race: that acceptance excludes an earlier accepted
update 1 at \(B\), and the strict update-number test rejects update 1
thereafter. Such evidence must identify the order and actual accepted
order update; a publisher-side send or unrelated progress is not
enough. Resets, competing issuers, and changes of order identity
would require separate evidence.

The pinned adapter tests a node's release flag before handing it to
the navigation callback. Its existing unreleased-node test exercises
that guard. Neither the new validator replay nor the guard certifies
a robot's physical clearance or the completion of an already
dispatched operation. A different action or source can still
occupy the bay, and that requires an independent check.

\paragraph{Comparator.}
This is an applicable stock mechanism to beat, not an improvement
over a planner given the same accepted anchor and update evidence.
It provides a real selective alternative for work still in the
horizon, with a boundary at the latest potentially accepted release.
No independent two-fault mission or I1 acceptance follows.
\end{document}
